\mode*
\only<presentation>{\setminted{autogobble,linenos,fontsize=\scriptsize}}

\section{A worked example: the patterns of variation in debugging}
\label{debugging}

Let us illustrate the argument of \cref{debugging-as-learning} with a worked
example: a debugging episode in which the patterns of variation appear one by
one.
We will use a particular feature related to default arguments in Python that
few expect, hence most readers will hopefully not know about this and get the
intended experience.
(Anyone who wants to learn about this aspect must go through a variant of
these patterns of variation, regardless of level; the phenomenon is simply
too advanced a topic for introductory programming, which is why we aim the
example at the instructors and researchers for whom this paper is intended.)

Here is the problem.
We assume familiarity with default values for arguments in Python.
The function \mintinline{python}{expand} below takes a list and appends an
element \mintinline{python}{1} to it; we call it a few times.

\begin{frame}[fragile]
  \begin{pyblock}[tryfirst]
def expand(x=[]):
  x.extend([1]) # x += [1]
  return x

print(
  f"[1] -> {expand([1])}\n"
  f"[2] -> {expand([2])}\n"
  f"()  -> {expand()}\n"
  f"()  -> {expand()}\n"
)
  \end{pyblock}
  \begin{block}{Output}\printpythontex[verbatim][highlightlines={4}]\end{block}

  \begin{exercise}\label{ex:tryfirst}
    The two last calls are identical, so they should reasonably print the
    same result---but they do not (highlighted line).
    Before reading on: debug this behaviour until you can explain it.
    While doing so, note \emph{how} you go about it---which examples you
    try, what you vary, and what you keep unchanged.
  \end{exercise}
\end{frame}
\ltnote{%
  Try-first design: the reader attempts the debugging---generates the
  variation themselves---before seeing the analysis, and the
  note-how-you-do-it prompt turns the reader's own episode into data to
  compare against the patterns below.
}

We now walk through the patterns of variation; compare them with the steps
you just took yourself.

\begin{description}
  \item[Contrast] The contrast pattern requires two examples to create
    contrast.
    The left-hand example below is the problem of \cref{ex:tryfirst}; the
    right-hand example is a well-behaved variant of the same function.
    The left-hand (familiar) example should be read first and then the
    right-hand example will highlight the changes made to the left-hand
    example to get the right-hand one.
    (So that one can make the changes oneself.)

    In an actual debugging session nobody serves us this contrast: the first
    part of the session is about \emph{finding} it.
    The debugger makes changes that introduce variation in different
    aspects, until hitting the aspect that produces the contrast---the
    critical one.
    That search depends on what aspects the debugger has already discerned.
    Here we skip that search and present the found contrast directly.

    \begin{frame}<presentation>
      Let's start with a \emph{contrasting} example \dots
    \end{frame}

    \begin{frame}[fragile]
    \begin{minipage}[t]{0.45\columnwidth}
      \begin{pyblock}[default2]
def expand(x=[]):
  x.extend([1]) # x += [1]
  return x

print(
  f"[1] -> {expand([1])}\n"
  f"[2] -> {expand([2])}\n"
  f"()  -> {expand()}\n"
  f"()  -> {expand()}\n"
)
      \end{pyblock}
      \begin{block}{Output}\printpythontex[verbatim][highlightlines={4}]\end{block}
    \end{minipage}
    \hfill
    \begin{minipage}[t]{0.45\columnwidth}
      \begin{pyblock}[default1][highlightlines={2}]
def expand(x=[]):
  return x + [1]


print(
  f"[1] -> {expand([1])}\n"
  f"[2] -> {expand([2])}\n"
  f"()  -> {expand()}\n"
  f"()  -> {expand()}\n"
)
      \end{pyblock}
      \begin{block}{Output}\printpythontex[verbatim][highlightlines={4}]\end{block}
    \end{minipage}
    \end{frame}

    On the right-hand side, the change (highlighted) is that we no longer
    modify \mintinline{python}{x}: we build and return a new list instead.
    In the output the weird behaviour is gone: the two identical calls now
    print the same result (highlighted).
    So the phenomenon seems connected to modifying the argument in place---it seems like the left-hand version appends on the value returned last
    time.

    Note that with the surface approach to learning one could stop here:
    \enquote{the problem is solved, remember to not make this mistake
    again.}
    Continuing to open up the phenomenon, as the following patterns do, is
    exactly the deep approach to learning.

    We continue the contrast with more calls.
    (For brevity we don't repeat the definition of the
    function~\mintinline{python}{expand}; we keep the two variants side by
    side, so the between-column comparison remains the contrast.)

    \begin{frame}[fragile]
    \begin{minipage}[t]{0.45\columnwidth}
      \begin{pyblock}[default2][highlightlines={4-7}]
print(
  f"()  -> {expand()}\n"
  f"[3] -> {expand([3])}\n"
  f"()  -> {expand()}\n"
  f"[]  -> {expand([])}\n"
  f"[]  -> {expand([])}\n"
  f"()  -> {expand()}\n"
)
      \end{pyblock}
      \begin{block}{Output}\printpythontex[verbatim][highlightlines={3-6}]\end{block}
    \end{minipage}
    \hfill
    \begin{minipage}[t]{0.45\columnwidth}
      \begin{pyblock}[default1]
print(
  f"()  -> {expand()}\n"
  f"[3] -> {expand([3])}\n"
  f"()  -> {expand()}\n"
  f"[]  -> {expand([])}\n"
  f"[]  -> {expand([])}\n"
  f"()  -> {expand()}\n"
)
      \end{pyblock}
      \begin{block}{Output}\printpythontex[verbatim]\end{block}
    \end{minipage}
    \end{frame}

    We can see in the output that when we use the default value in the
    argument, the return value keeps expanding.
    These calls contrast along several different aspects.
    Between the columns we still contrast the function variant: modifying
    the default in place against building a new list.
    Within each column we contrast how the argument is supplied: the
    highlighted calls show that passing the empty list explicitly,
    \mintinline{python}{expand([])}, always gives \mintinline{python}{[1]},
    while using the default, \mintinline{python}{expand()}, keeps growing.
    And among the explicit arguments we contrast an arbitrary list,
    \mintinline{python}{[3]}, against the empty list,
    \mintinline{python}{[]}---the same value as the default.
    In the highlighted calls, what varies is only \emph{how} the empty list
    is supplied---so the phenomenon is tied to the default object itself,
    not to empty lists in general.
    This gives us a hypothesis about which aspect is the critical one: the
    phenomenon appears because we modify the default argument in place, and
    because the default argument is an object, which Python handles by
    reference.

  \item[Generalisation] This brings us to the generalisation pattern.
    In the generalisation pattern we vary the non-critical aspects and keep the
    critical aspect invariant.
    In a sense, the generalisation pattern is the same as induction or the
    scientific method.

    \begin{frame}<presentation>
      \dots and then \emph{generalise} it.
    \end{frame}

    \begin{frame}<presentation>
      \begin{remark}[Generalisation]
        \begin{itemize}
          \item We vary the non-critical aspects.
          \item We keep the critical aspect invariant.
            \pause
          \item \alert<3>{This is also known as the Scientific Method.}
        \end{itemize}
      \end{remark}
    \end{frame}

    \begin{frame}<presentation>[fragile]
      \begin{example}[experiment2.py]
        \begin{onlyenv}<1>
          \inputminted[firstline=1,lastline=13]{python}{examples/experiment2.py}
        \end{onlyenv}
        \begin{onlyenv}<2>
          \inputminted[firstline=1,lastline=2]{python}{examples/experiment2.py}
          \inputminted[firstline=9,lastline=10]{python}{examples/experiment2.py}
          \inputminted[firstline=14]{python}{examples/experiment2.py}
        \end{onlyenv}
      \end{example}
    \end{frame}

    \begin{frame}<presentation>[fragile]
      \begin{example}[experiment3.py]
      \inputminted{python}{examples/experiment3.py}
      \end{example}
    \end{frame}

    Now, we'll change the example, but keep this property invariant---we
    still modify the default argument, an object handled by reference;
    or, phrased in terms of the scientific method, we try to falsify our
    hypothesis.
    (We highlight the lines that keep this property, \ie remain invariant in
    the variation theoretic sense.)
    \begin{frame}[fragile]
    \begin{pyblock}[default1][highlightlines={10-12}]
class Person:
  def __init__(self, first, last):
    self.first = first
    self.last = last

  def __str__(self):
    return f"{self.first} {self.last}"

def create_person(first=None, last=None,
                  person_base=Person("Gina", "Jones")):
  if first: person_base.first = first
  if last: person_base.last = last
  return person_base

person_default = create_person()
print(person_default)
person_A = create_person("Ada")
print(person_A)
person_B = create_person("Beda")
print(person_B)
print(person_A)
print(person_default)
    \end{pyblock}
    \begin{block}{Output}\printpythontex[verbatim][highlightlines={4-5}]\end{block}
    \end{frame}

    We can see that the last three lines of the output is the same.
    The first of them is expected, but we expected the reprint of
    \mintinline{python}{person_A} to still give \enquote{Ada Jones} and
    \mintinline{python}{person_default} to still give
    \enquote{Gina Jones}---the default.
    Instead all three print \enquote{Beda Jones} (highlighted), indicating
    that \mintinline{python}{person_A}, \mintinline{python}{person_B} and
    \mintinline{python}{person_default} all seem to refer to the same
    object.

  \item[Fusion] Now we can fuse this back with our previous understanding of
    default arguments.
    First we test that it doesn't work with non-mutables, like integers.
    \begin{frame}<presentation>
      \dots and then \emph{fuse} it back.
    \end{frame}
    \begin{frame}[fragile]
    \begin{pyblock}[default1]
def increment(x=1):
  x += 1
  return x

print(f"(1) -> {increment(1)}")
print(f"(2) -> {increment(2)}")
print(f"()  -> {increment()}")
print(f"()  -> {increment()}")
print(f"()  -> {increment()}")
    \end{pyblock}
    \begin{block}{Output}\printpythontex[verbatim][highlightlines={3-5}]\end{block}
    \end{frame}

    And we can thus conclude that this phenomenon happens only with mutable
    objects that are handled by reference in Python.

    This must be due to how Python is constructed.
    The default value seems to be constructed when the function is defined, not
    when it's called (like in other languages, C++ for instance), and then
    referenced (not copied) whenever the function is called without the
    argument.
    Consider this example.

    \begin{frame}[fragile]
    \begin{pyblock}[default1][highlightlines=5]
class TraceClass:
  def __init__(self):
    print(f"{self} created")

def test_function(obj=TraceClass()):
  print(f"test_function called with obj = {obj}")

print("Test code begins")
test_function()
    \end{pyblock}
    \begin{block}{Output}\printpythontex[verbatim][highlightlines={1}]\end{block}
    \end{frame}

    We see that the print statement from the constructor is executed before the
    test code is executed (highlighted line), supporting our hypothesis that
    the default value is constructed when the function is defined, not the when
    function is called, and then referenced throughout.
    This example fuses all the aspects: the construction time of the default
    value, that it is referenced rather than copied, and that modifying it in
    place is what makes the phenomenon visible.

    We can also use an example from another language: C++.
    In C++, default arguments are evaluated when the function is called, not
    when it's defined.
    \begin{frame}[fragile]
      \inputminted[firstline=4]{cpp}{examples/experiment5.cpp}
      \begin{block}{Output}
        \inputminted[linenos=false,highlightlines={2,4}]{text}{examples/experiment5.out}
      \end{block}
    \end{frame}

    We see that in C++ the object is created anew at every call, after the
    test code has started (highlighted lines)---the opposite of Python.
    (To mirror Python's reference semantics we pass the argument by
    reference; C++ does not allow binding a plain reference,
    \mintinline{cpp}{TraceClass&}, to the default temporary, so it must be
    an rvalue reference, \mintinline{cpp}{TraceClass&&}.
    Even so, a new object is constructed at each call.)
    Note that the two printed addresses may even coincide; that does not
    mean it is the same object, only that the same memory location is
    reused for the two calls' temporaries---the constructor runs anew each
    call.
    That the objects are different is also supported by the modification:
    \mintinline{cpp}{test_function} appends an element to the object, yet
    the size is \(1\) in both calls---the modification of the first call's
    default does not carry over to the second.
\end{description}

By now, we fully understand this phenomenon, and we see default values in a
more powerful way.

\subsection{Consequences}

There are several things to note with this example.
First, note how the contrast pattern is designed to focus your (the
reader's) attention to the phenomenon at hand: that default values can change
during execution.
Next, the generalisation pattern broadens our view of when this phenomenon
happens, that the objects are referenced and reused.
Finally, the fusion pattern merges this back into our original view of default
values, namely that they work as usual for non-mutable types (\eg integers).
By now we completely understand the phenomenon at hand: the patterns did not
merely remove the bug, they produced an explanation.

In the terms of \citeauthor{NCOL}'s account of discovery
(\cref{debugging-as-learning}), what we did with this example is to open up
the dimension of construction/execution time for default arguments: once
that dimension is opened, construction at definition time and construction
at call time become two features that can vary---and the C++ comparison is
exactly that variation.

Second thing to note, if you are a seasoned programmer, once you had that
initial contrast pattern the remaining patterns (generalisation and fusion)
probably resembles quite a lot what you would have tested yourself to make
sense of this phenomenon---it would probably resemble how you would go about to
debug this and explain it to yourself.
Unless, of course, you opted for the surface approach to learning.

But then the main issue of debugging is finding that initial contrast.
Why, then, would someone using the deep approach to learning do better?
Because the search for the contrast is itself a matter of variation: the
debugger varies aspect after aspect until the critical one responds, and
which aspects one can vary at all depends on what one has already
discerned.
In our example, to even conceive of the variations above, the debugger must
already see---in a more powerful way---at least the following aspects.
\begin{enumerate}
  \item That a parameter may get its value from a default, so that
    supplying the value explicitly is a variation at all.
  \item That a list is a mutable object, which can be modified in place or
    rebuilt.
  \item That variables and parameters hold references to objects, so that
    two names may refer to the same object.
  \item That objects are constructed, and that construction happens through
    execution.
\end{enumerate}
These are basically the prerequisites for the search: an aspect one cannot
yet discern is an aspect one cannot vary.
Note that \emph{when} the default argument's definition is executed---and
hence when the default is constructed---is not among the prerequisites:
that is precisely the dimension the search opens up.
And note that it is \emph{the bug} that opens the dimension, by contrasting the
actual behaviour with what was expected; the search then locates which
aspect carries that contrast.
The deep approach is exactly the disposition to generate that variation;
the surface approach offers no search procedure at all---it can only
recognise mistakes already made and memorised.

\subsection{A pilot with experienced programmers}
\label{pilot}

We actually tested this hypothesis on colleagues who have been
programming for many years, are well-versed in Python but didn't know of this
phenomenon.
We gave them the same problem as in \cref{ex:tryfirst} and asked them to
debug this behaviour and later explain it when they understood it.
(To record the data, we asked them to think aloud and recorded their screen and
voice in a Zoom session.)
Nowadays we can also record such sessions directly with learnlog
\autocite{learnlog}, as we do in the study in \cref{method}.

% XXX Test with more colleagues.
% Tested with Alexander.
% TODO Expand this section with details from the recorded sessions (see the
% KTH Play playlist referenced in debug_data/data.md): which concrete
% examples each colleague tried, in what order, and how the patterns map
% onto them.  The companion paper only sketches this observation; the full
% account belongs here.
The concrete examples that they tried varied from person to person, some tried
many more examples than above, but the patterns of variation shown above were
present.
Indeed, they couldn't explain the phenomenon until they had generated all the
patterns above---covering both generalisation and fusion.
Their behaviour also matches what the expert debugging literature
prescribes: \textcite{WhyProgramsFail} teaches debugging as the
scientific method---creating and verifying hypotheses about the failure
through experiments---which in variation-theoretic terms is generating
the contrast and generalisation patterns.

So far we have two recorded sessions.
We transcribe the recordings using learnlog---replaying the sessions'
edit--run cycles into learnlog logs---so that the transcripts can be
analysed with the episode coding of \cref{method}, and the coded episodes
reported here.
% TODO Transcribe and analyse the two recorded sessions (KTH Play playlist in
% debug_data/data.md); report the coded episodes here.
% XXX Stub: per-session summaries to be added.
% - Session 1: examples tried, in what order, which patterns they realise.
% - Session 2: same.

\mode<presentation>{%
\begin{frame}
  \begin{remark}[Tested]
    \begin{itemize}
      \item Tested this on a colleague.
      \item He couldn't explain the phenomenon until he had generated all the
        patterns above.
      \item Plan to do this with more colleagues.
    \end{itemize}
  \end{remark}
\end{frame}%
}

This pilot supports \cref{h:patterns} for experienced programmers: the
successful debugging episodes contained all the patterns of variation, and the
explanation did not come until the patterns were complete.
What the pilot cannot tell us is how this plays out for \emph{students}, nor
whether the ability to generate the patterns is connected to the deep approach
to learning.
That is the purpose of the study we describe next.
